\documentclass[aspectratio=169]{beamer}

\usepackage[utf8]{inputenc}
\usepackage[T1]{fontenc}
\usepackage[brazil]{babel}
\usepackage{amsmath, amssymb}
\usepackage{microtype}
\usepackage{csquotes}
\usepackage{natbib}
\usepackage{tikz}

% Colors
\definecolor{mblack}{RGB}{20, 20, 20}
\definecolor{mprimary}{RGB}{50, 98, 115}     % dark blue
\definecolor{maccent}{RGB}{227, 151, 116}    % orange
\definecolor{mwhite}{RGB}{255, 255, 255}

% Theme
\usetheme{default}
\usecolortheme{default}
\setbeamercolor{normal text}{fg=mblack, bg=white}
\setbeamercolor{frametitle}{fg=white, bg=mprimary}
\setbeamercolor{title}{fg=white}
\setbeamercolor{author}{fg=white}
\setbeamercolor{date}{fg=white}
\setbeamercolor{itemize item}{fg=maccent}
\setbeamercolor{itemize subitem}{fg=mprimary}
\setbeamercolor{enumerate item}{fg=maccent}
\setbeamercolor{enumerate subitem}{fg=mprimary}
\setbeamercolor{block title}{fg=white, bg=mprimary}
\setbeamercolor{block body}{fg=mblack, bg=mprimary!8}
\setbeamercolor{block title alerted}{fg=white, bg=maccent}
\setbeamercolor{block body alerted}{fg=mblack, bg=maccent!10}
\setbeamercolor{alerted text}{fg=maccent}

\setbeamerfont{frametitle}{series=\bfseries, size=\large}
\setbeamerfont{title}{series=\bfseries, size=\LARGE}
\setbeamerfont{block title}{series=\bfseries}

\setbeamertemplate{navigation symbols}{}

% Footer
\setbeamertemplate{footline}{%
  \leavevmode%
  \hbox{%
    \begin{beamercolorbox}[wd=0.7\paperwidth, ht=2.25ex, dp=1ex, leftskip=0.3cm]{normal text}%
      \tiny\color{mblack!50} Filosofia da Ciência -- Falsificacionismo
    \end{beamercolorbox}%
    \begin{beamercolorbox}[wd=0.3\paperwidth, ht=2.25ex, dp=1ex, rightskip=0.3cm, right]{normal text}%
      \tiny\color{mblack!50}\insertframenumber{} / \inserttotalframenumber
    \end{beamercolorbox}%
  }%
}

% Frametitle with accent rule
\setbeamertemplate{frametitle}{%
  \nointerlineskip%
  \begin{beamercolorbox}[wd=\paperwidth, ht=0.55cm, dp=0.2cm, leftskip=0.5cm]{frametitle}%
    \usebeamerfont{frametitle}\insertframetitle%
  \end{beamercolorbox}%
  \vspace{-0.1cm}%
  \textcolor{maccent}{\rule{\paperwidth}{2pt}}%
}

% Title page: three-band layout
\setbeamercolor{titlebar top}{bg=mwhite}
\setbeamercolor{titlebar main}{bg=mprimary}
\setbeamercolor{titlebar bottom}{bg=maccent}

\setbeamertemplate{title page}{%
  \vspace{0pt}%
  \begin{beamercolorbox}[wd=\paperwidth, ht=0.55cm, dp=0pt]{titlebar top}\end{beamercolorbox}%
  \begin{beamercolorbox}[wd=\paperwidth, ht=3.8cm, dp=0pt, center]{titlebar main}%
    \vspace{0.6cm}
    {\usebeamerfont{title}\color{white}\inserttitle\par}%
    \vspace{0.35cm}{\color{white}\small\insertauthor\par}%
    \vspace{0.15cm}{\color{white}\footnotesize\insertdate\par}%
    \vspace{0.25cm}
  \end{beamercolorbox}%
  \begin{beamercolorbox}[wd=\paperwidth, ht=1.3cm, dp=0pt]{titlebar bottom}\end{beamercolorbox}%
}

\setlength{\itemsep}{4pt}

% Highlight commands
\newcommand{\impt}[1]{\textcolor{maccent}{\textbf{#1}}}
\newcommand{\ntc}[1]{\textcolor{mprimary}{#1}}

\title{Falsificacionismo}
\author{Filosofia da Ciência\footnote{Baseado em \citet{chalmers1993}, caps.~IV, V e VI.}}
\date{\today}

\begin{document}

\begin{frame}[plain]\titlepage\end{frame}

\begin{frame}{Sumário}\tableofcontents\end{frame}

% =====================================================
\section{Motivação}
% =====================================================

\begin{frame}{Por que o falsificacionismo?}
O \impt{indutivismo} enfrentava dois problemas graves:
\vspace{0.3cm}
\begin{itemize}
  \item Problema da \impt{indução} --- sem justificação lógica
  \item Observação \impt{não é neutra} --- é carregada de teoria
\end{itemize}
\pause
\vspace{0.5cm}
\begin{block}{A proposta falsificacionista}
Abandonar a pretensão de \emph{provar} teorias. \\
A ciência avança por \impt{conjecturas ousadas} que a experiência pode \impt{refutar}.
\end{block}
\pause
\vspace{0.3cm}
\ntc{Principal representante:} Karl Popper.
\end{frame}

\begin{frame}{A ideia central}
\begin{alertblock}{}
A ciência não é um corpo de verdades provadas. \\
É um conjunto de hipóteses \impt{conjecturais} submetidas a testes cada vez mais severos.
\end{alertblock}
\vspace{0.5cm}
\pause
Teorias não são \emph{confirmadas} --- no máximo, \impt{corroboradas}. \\
Teorias podem ser \emph{refutadas}.
\end{frame}

% =====================================================
\section{Falsificacionismo ingênuo (Cap.~IV)}
% =====================================================

\begin{frame}{A assimetria lógica}
Considere: \emph{``Todos os corvos são pretos.''}
\vspace{0.4cm}
\pause
\begin{columns}[t]
\column{0.48\textwidth}
\begin{block}{Verificação}
Nenhum número finito de corvos pretos \impt{prova} a lei. \\[0.2cm]
O próximo pode ser branco.
\end{block}

\column{0.48\textwidth}
\pause
\begin{block}{Falsificação}
\impt{Um único} corvo não-preto \\
\impt{refuta} a lei definitivamente.
\end{block}
\end{columns}
\vspace{0.5cm}
\pause
\ntc{Lógica por trás:}
\[
(T \rightarrow O) \wedge \neg O \;\vdash\; \neg T \qquad \text{(modus tollens)}
\]
Já $(T \rightarrow O) \wedge O$ \emph{não} implica $T$ (falácia de afirmar o consequente).
\end{frame}

\begin{frame}{Falseabilidade como critério de demarcação}
\begin{block}{Definição}
Uma teoria é \impt{falsificável} se existe ao menos um enunciado observacional logicamente possível que é \emph{incompatível} com ela.
\end{block}
\pause
\vspace{0.4cm}
\ntc{Papel:} separa ciência de não-ciência.
\pause
\vspace{0.3cm}
\begin{itemize}
  \item Falsificável $\neq$ falso
  \item Falsificável $=$ arrisca-se a estar errado
\end{itemize}
\end{frame}

\begin{frame}{Exemplos --- falsificáveis}
\begin{itemize}
  \item ``Nunca chove às quartas-feiras.''
  \pause
  \item ``Todos os corpos em queda livre caem com $a \approx 9{,}8\,\mathrm{m/s^2}$.''
  \pause
  \item ``Os planetas descrevem órbitas elípticas em torno do Sol.''
\end{itemize}
\vspace{0.5cm}
\pause
\ntc{Cada uma proíbe algo.} \\
Cada uma pode ser derrubada por uma observação.
\end{frame}

\begin{frame}{Exemplos --- não-falsificáveis}
\begin{itemize}
  \item ``Ou chove ou não chove.'' \hfill (tautologia)
  \pause
  \item ``Todos os pontos do círculo são equidistantes do centro.'' \hfill (definição)
  \pause
  \item ``Pode ou não chover amanhã.'' \hfill (compatível com tudo)
\end{itemize}
\vspace{0.5cm}
\pause
\begin{alertblock}{}
Podem ser verdadeiras --- mas \impt{não dizem nada sobre o mundo}.
\end{alertblock}
\end{frame}

\begin{frame}{A crítica às teorias ``explicam-tudo''}
Popper dirige essa crítica a:
\vspace{0.3cm}
\begin{itemize}
  \item Psicanálise freudiana
  \item Certas formulações do marxismo
  \item Astrologia
\end{itemize}
\pause
\vspace{0.5cm}
\begin{alertblock}{O defeito}
Não é serem falsas. \\
É serem \impt{compatíveis com qualquer observação}.
\end{alertblock}
\pause
\vspace{0.3cm}
\ntc{Uma teoria que não proíbe nada, não diz nada.}
\end{frame}

\begin{frame}{Graus de falseabilidade}
Uma boa teoria não basta ser falsificável --- deve ser \impt{muito falsificável}.
\vspace{0.4cm}
\pause
\begin{block}{Exemplo}
\begin{itemize}
  \item $T_1$: Marte move-se em órbita fechada ao redor do Sol.
  \item $T_2$: Todos os planetas movem-se em órbitas fechadas.
  \item $T_3$: Todos os planetas movem-se em órbitas \impt{elípticas}.
\end{itemize}
\end{block}
\pause
\vspace{0.3cm}
\ntc{Ordem de falseabilidade:} \quad $T_3 > T_2 > T_1$
\pause
\vspace{0.3cm}
\begin{itemize}
  \item[$\rightarrow$] O falsificacionista prefere a teoria mais \impt{ousada}.
\end{itemize}
\end{frame}

\begin{frame}{Clareza e precisão}
Quanto mais \impt{precisa} a teoria, mais falsificável.
\vspace{0.5cm}
\pause
\begin{columns}[t]
\column{0.48\textwidth}
\begin{block}{Vago}
``A Lua é mais ou menos esférica.''
\end{block}

\column{0.48\textwidth}
\pause
\begin{block}{Preciso}
``A Lua é uma esfera de raio $1737{,}4$\,km.''
\end{block}
\end{columns}
\vspace{0.5cm}
\pause
\begin{alertblock}{}
Vagueza = \impt{escudo contra refutação}. \\
Teorias vagas reinterpretam-se \emph{ad hoc} diante de qualquer resultado.
\end{alertblock}
\end{frame}

\begin{frame}{Falsificacionismo e progresso (ingênuo)}
A história da ciência seria uma sucessão de:
\vspace{0.3cm}
\begin{enumerate}
  \item Conjecturas ousadas
  \pause
  \item Refutações
  \pause
  \item Conjecturas ainda mais ousadas
\end{enumerate}
\pause
\vspace{0.5cm}
\begin{block}{Progresso científico}
Não há acúmulo de verdades provadas. \\
Há \impt{eliminação de erros}.
\end{block}
\end{frame}
% =====================================================
\begin{frame}{O problema dos morcegos (Chalmers, Cap.~IV)}
\ntc{Observação inicial:} 
\vspace{0.2cm}
\begin{block}{}
\centering Morcegos voam com agilidade e capturam insetos \\ mesmo em condições de muito pouca luz.
\end{block}
\pause
\vspace{0.3cm}
\ntc{Problema:} \impt{como?}
\pause
\vspace{0.4cm}
\begin{alertblock}{O falsificacionista em ação}
Propõe \impt{conjecturas ousadas} e submete cada uma a \impt{testes severos}. \\
Conjecturas refutadas são \emph{eliminadas}; a sobrevivente fica \impt{corroborada}.
\end{alertblock}
\end{frame}

\begin{frame}{Conjectura 1 --- Visão}
\begin{block}{Conjectura}
Morcegos enxergam normalmente, mesmo com luz muito fraca.
\end{block}
\pause
\vspace{0.3cm}
\ntc{Predição arriscada:} morcegos com os olhos vendados \impt{devem} colidir com obstáculos e falhar em capturar presas.
\pause
\vspace{0.4cm}
\begin{alertblock}{Teste}
Vendam-se os olhos dos morcegos.\\
\impt{Resultado:} continuam voando e caçando normalmente.
\end{alertblock}
\pause
\vspace{0.3cm}
\ntc{Conjectura 1 está \impt{falsificada}.}
\end{frame}

\begin{frame}{Conjectura 2 --- Olfato}
\begin{block}{Conjectura}
Morcegos detectam obstáculos e presas pelo olfato.
\end{block}
\pause
\vspace{0.3cm}
\ntc{Predição arriscada:} morcegos com o olfato bloqueado devem colidir e falhar em caçar.
\pause
\vspace{0.4cm}
\begin{alertblock}{Teste}
Bloqueia-se o nariz dos morcegos.\\
\impt{Resultado:} continuam voando e caçando normalmente.
\end{alertblock}
\pause
\vspace{0.3cm}
\ntc{Conjectura 2 está \impt{falsificada}.}
\end{frame}

\begin{frame}{Conjectura 3 --- Audição}
\begin{block}{Conjectura}
Morcegos orientam-se pela audição, detectando ecos de sons por eles próprios emitidos.
\end{block}
\pause
\vspace{0.3cm}
\ntc{Predições arriscadas:}
\begin{itemize}
  \item Morcegos com \impt{ouvidos tapados} devem perder a orientação.
  \item Morcegos com \impt{boca fechada} (sem emitir sons) também devem perder.
\end{itemize}
\pause
\vspace{0.4cm}
\begin{alertblock}{Teste}
Ambas as predições são \impt{confirmadas}.
\end{alertblock}
\pause
\vspace{0.3cm}
\ntc{Conjectura 3 sobrevive e fica \impt{corroborada}.}
\end{frame}

\begin{frame}{O que esse exemplo ilustra}
\begin{itemize}
  \item A ciência procede por \impt{tentativa e eliminação de erros}
  \pause
  \item Boas conjecturas são \impt{falsificáveis} --- arriscam-se a estar erradas
  \pause
  \item Testes severos \impt{eliminam} as falsas
  \pause
  \item A conjectura sobrevivente \impt{não está provada} --- está \emph{corroborada}
\end{itemize}
\pause
\vspace{0.4cm}
\begin{alertblock}{}
Amanhã, uma nova observação poderia refutar a hipótese da ecolocalização. \\
\impt{Toda teoria permanece conjectural.}
\end{alertblock}
\end{frame}

% =====================================================
\begin{frame}{De Aristóteles a Einstein}
\ntc{Uma sequência de teorias cada vez mais ousadas:}
\vspace{0.4cm}
\begin{columns}[t]
\column{0.32\textwidth}
\begin{block}{Aristóteles}
Corpos caem porque buscam seu ``lugar natural''.\\[0.2cm]
\impt{Física qualitativa.}
\end{block}

\column{0.32\textwidth}
\pause
\begin{block}{Newton}
Lei universal da gravitação.\\[0.2cm]
\impt{Quantitativa e preditiva.}\\
Explica Kepler e Galileu.
\end{block}

\column{0.32\textwidth}
\pause
\begin{block}{Einstein}
Gravidade = curvatura do espaço-tempo.\\[0.2cm]
\impt{Novas predições:} curvatura da luz, precessão de Mercúrio.
\end{block}
\end{columns}
\pause
\vspace{0.5cm}
\begin{alertblock}{Progresso no sentido falsificacionista}
Cada teoria é \impt{mais falsificável} que a anterior, \impt{explica seu sucesso} e faz \impt{novas predições corroboradas}.
\end{alertblock}
\end{frame}

% =====================================================
\section{Falsificacionismo sofisticado (Cap.~V)}
% =====================================================

\begin{frame}{Problema do falsificacionismo ingênuo}
Muitas teorias importantes \impt{nasceram refutadas}.
\vspace{0.5cm}
\pause
\begin{block}{Newton e a órbita da Lua}
Primeiro cálculo de Newton discordava dos dados. \\
Um falsificacionista rigoroso teria abandonado a gravitação. \\
\pause
\impt{Erro estava em um dado observacional} (diâmetro da Terra).
\end{block}
\pause
\vspace{0.4cm}
\ntc{Lição:} abandonar teorias ao primeiro tropeço seria desastroso.
\end{frame}

\begin{frame}{A reformulação sofisticada}
\ntc{Deslocamento de foco:} de teorias isoladas para \impt{sequências de teorias}.
\vspace{0.4cm}
\pause
\begin{block}{Critério sofisticado}
$T'$ substitui legitimamente $T$ quando:
\begin{enumerate}
  \item $T'$ é \impt{mais falsificável} que $T$
  \item $T'$ explica o sucesso prévio de $T$
  \item $T'$ faz \impt{novas predições} --- e parte delas é \impt{corroborada}
\end{enumerate}
\end{block}
\pause
\vspace{0.3cm}
\ntc{Formulação central:} Imre Lakatos.
\end{frame}
%%%%%%%%%%%%%%%%%%%%%%%%%%%%%%%%%%%%%%%%%%
%%%%%%%%%%%%%%%%%%%%%%%%%%%%%%%%%%%%%%%%%%
\begin{frame}{O problema: o que fazer diante de uma refutação?}
Uma teoria $T$ prediz $O$, mas observa-se $\neg O$.
\vspace{0.4cm}
\pause
\ntc{Duas saídas possíveis:}
\vspace{0.3cm}
\begin{enumerate}
  \item \impt{Abandonar} $T$ e buscar teoria melhor
  \pause
  \item \impt{Modificar} $T$ para acomodar $\neg O$
\end{enumerate}
\pause
\vspace{0.5cm}
\begin{alertblock}{A pergunta central}
Quando a modificação (saída 2) é \impt{legítima}? \\
Quando ela é apenas um \impt{truque} para salvar a teoria?
\end{alertblock}
\end{frame}

\begin{frame}{Modificações \emph{ad hoc}}
\begin{block}{Definição}
Uma modificação é \impt{\emph{ad hoc}} quando é feita \impt{apenas} para salvar a teoria da refutação, \impt{sem gerar nenhuma nova consequência testável}.
\end{block}
\pause
\vspace{0.5cm}
\ntc{Do latim \emph{ad hoc}:} ``para isto'', ``feita sob medida''.
\pause
\vspace{0.3cm}
\begin{itemize}
  \item A modificação resolve \impt{só aquele} contraexemplo
  \item Não diz nada sobre outros casos
  \item \impt{Não arrisca} nenhuma predição nova
\end{itemize}
\end{frame}

\begin{frame}{Exemplo ingênuo de modificação \emph{ad hoc}}
\ntc{Teoria original:} ``Todos os corvos são pretos.''
\pause
\vspace{0.3cm}
\ntc{Observação:} um corvo branco.
\pause
\vspace{0.5cm}
\begin{alertblock}{Modificação \emph{ad hoc}}
``Todos os corvos são pretos, \impt{exceto este}.''
\end{alertblock}
\pause
\vspace{0.4cm}
\ntc{Por que é ilegítima?}
\begin{itemize}
  \item Salva a teoria do contraexemplo --- mas \impt{só desse}
  \pause
  \item \impt{Não prediz} nada novo
  \pause
  \item \impt{Não é testável} em nenhum caso além do que motivou a modificação
\end{itemize}
\end{frame}

\begin{frame}{Exemplo de modificação \impt{legítima}}
\ntc{Teoria:} mecânica newtoniana aplicada ao sistema solar.
\pause
\vspace{0.3cm}
\ntc{Problema (1840s):} órbita de Urano não bate com as predições.
\pause
\vspace{0.4cm}
\begin{block}{Modificação proposta (Le Verrier, Adams)}
Existe um \impt{planeta ainda não descoberto} perturbando Urano.
\end{block}
\pause
\vspace{0.3cm}
\ntc{Por que \emph{não} é \emph{ad hoc}?}
\begin{itemize}
  \item Faz \impt{predições novas e precisas}: massa, órbita, posição no céu
  \pause
  \item \impt{É testável} --- basta apontar o telescópio
  \pause
  \item \impt{Aumenta o conteúdo empírico} da teoria
\end{itemize}
\pause
\vspace{0.3cm}
\begin{alertblock}{}
\impt{Netuno foi observado em 1846}, na posição prevista.
\end{alertblock}
\end{frame}

\begin{frame}{A diferença em uma linha}
\begin{columns}[t]
\column{0.48\textwidth}
\begin{alertblock}{\emph{Ad hoc}}
Fecha um buraco. \\[0.2cm]
\impt{Não abre nenhuma janela.}
\end{alertblock}

\column{0.48\textwidth}
\pause
\begin{block}{Legítima}
Fecha um buraco. \\[0.2cm]
\impt{Abre novas janelas} --- novos testes possíveis.
\end{block}
\end{columns}
\pause
\vspace{0.6cm}
\begin{block}{Critério}
Uma modificação é legítima sse \impt{aumenta o conteúdo empírico} da teoria --- ou seja, se faz \impt{novas predições testáveis}.
\end{block}
\end{frame}

\begin{frame}{Novas predições e crescimento da ciência}
\ntc{Generalizando o ponto:} a ciência progride quando novas teorias (ou modificações) fazem \impt{predições novas e arriscadas} que depois são corroboradas.
\pause
\vspace{0.4cm}
\begin{block}{Critério de progresso (falsificacionismo sofisticado)}
$T'$ substitui legitimamente $T$ quando:
\begin{enumerate}
  \item $T'$ é \impt{mais falsificável} que $T$
  \item $T'$ explica o sucesso prévio de $T$
  \item $T'$ faz \impt{novas predições} --- e parte delas é \impt{corroborada}
\end{enumerate}
\end{block}
\end{frame}

\begin{frame}{Predições ousadas na história da ciência}
\begin{itemize}
  \item \impt{Maxwell} (1865): ondas eletromagnéticas à velocidade da luz \\
        \ntc{Confirmado:} Hertz, 1887
  \pause
  \vspace{0.2cm}
  \item \impt{Einstein} (1915): curvatura da luz pelo Sol \\
        \ntc{Confirmado:} eclipse de 1919 (Eddington)
  \pause
  \vspace{0.2cm}
  \item \impt{Le Verrier} (1846): novo planeta além de Urano \\
        \ntc{Confirmado:} descoberta de Netuno no mesmo ano
\end{itemize}
\pause
\vspace{0.4cm}
\begin{alertblock}{}
Em todos os casos: a teoria predisse \impt{antes} de o fenômeno ser conhecido.\\
Isso é \impt{corroboração genuína} --- não mera acomodação.
\end{alertblock}
\end{frame}

\begin{frame}{Corroboração}
\begin{block}{Definição}
Uma teoria está \impt{corroborada} quando sobreviveu a testes severos, especialmente quando suas predições ousadas foram confirmadas.
\end{block}
\pause
\vspace{0.5cm}
\begin{alertblock}{Atenção}
Corroboração \impt{não é prova}. \\
Reflete o \emph{passado} da teoria --- não sua probabilidade de ser verdadeira.
\end{alertblock}
\pause
\vspace{0.4cm}
\ntc{Popper resiste:} corroboração \impt{não} deve ser lida como evidência indutiva. \\
Teorias altamente corroboradas podem ser refutadas amanhã.
\end{frame}
%%%%%%%%%%%%%%%%%%%%%%%%%%%%%%%%%%%%%%%%%%
%%%%%%%%%%%%%%%%%%%%%%%%%%%%%%%%%%%%%%%%%%
% =====================================================
\section{Limitações do falsificacionismo (Cap.~VI)}
% =====================================================

\begin{frame}{Três grandes limitações}
Mesmo na versão sofisticada, o falsificacionismo falha:
\vspace{0.4cm}
\pause
\begin{enumerate}
  \item \impt{Dependência} da observação em relação à teoria
  \pause
  \vspace{0.2cm}
  \item Tese de \impt{Duhem--Quine} (holismo da confirmação)
  \pause
  \vspace{0.2cm}
  \item Inadequação \impt{histórica}
\end{enumerate}
\end{frame}

\begin{frame}{1. Observação carregada de teoria}
Não existe observação \impt{neutra}. \\
Todo enunciado observacional $O$ pressupõe alguma teoria $T^*$.
\vspace{0.5cm}
\pause
\begin{block}{Consequência}
O embate ``$T$ vs $O$'' não é teoria vs fato bruto. \\
É um \impt{embate entre teorias}.
\end{block}
\pause
\vspace{0.4cm}
\ntc{Colapso:} enunciados observacionais são \impt{falíveis}. \\
A falsificação perde seu \impt{privilégio lógico}.
\end{frame}

\begin{frame}{O que é o ``privilégio lógico'' da falsificação?}
\vspace{0.3cm}
\[
(T \rightarrow O) \wedge \neg O \;\vdash\; \neg T \qquad \text{(modus tollens)}
\]
\pause
\vspace{0.3cm}
\begin{block}{O privilégio}
Se $\neg O$ é \impt{certo}, então $\neg T$ é \impt{certo}. \\
A falsificação seria \impt{logicamente conclusiva}. \\
A verificação nunca é.
\end{block}
\pause
\vspace{0.4cm}
\ntc{Toda a força do falsificacionismo} depende desse privilégio.
\end{frame}

\begin{frame}{Por que o privilégio se perde}
O \emph{modus tollens} entrega certeza sobre $\neg T$ \impt{apenas} na medida em que há certeza sobre $\neg O$.
\pause
\vspace{0.4cm}
\begin{alertblock}{Mas\ldots}
Enunciados observacionais \impt{pressupõem teorias} $T^*$ \\
(instrumentos, percepção, fenômenos de fundo).\\[0.2cm]
Logo, $\neg O$ é ele mesmo \impt{falível}.
\end{alertblock}
\pause
\vspace{0.4cm}
\ntc{Consequência:}
\[
(T \rightarrow O) \wedge \underbrace{\neg O}_{\text{falível}} \;\vdash\; \underbrace{\neg T}_{\text{igualmente falível}}
\]
\pause
\vspace{0.2cm}
A falsificação é agora \impt{tão segura quanto} o enunciado observacional que a sustenta --- e este é \impt{revisável}.
\end{frame}

\begin{frame}{De confronto teoria--fato a confronto teoria--teoria}
\begin{columns}[t]
\column{0.48\textwidth}
\begin{block}{Visão ingênua}
$T$ \emph{vs.} \impt{fato bruto $\neg O$} \\[0.2cm]
A falsificação é conclusiva.
\end{block}

\column{0.48\textwidth}
\pause
\begin{alertblock}{Realidade}
$T$ \emph{vs.} $T^*$ \\[0.2cm]
Duas teorias concorrendo. \\
\impt{Nenhuma com prioridade lógica.}
\end{alertblock}
\end{columns}
\pause
\vspace{0.5cm}
\ntc{É sempre possível} preservar $T$ rejeitando $T^*$, ou vice-versa. \\
A \impt{lógica não decide} qual abandonar.
\end{frame}

\begin{frame}{2. A tese de Duhem--Quine}
Teorias \impt{nunca} são testadas isoladamente.
\vspace{0.3cm}
\pause
\[
T \;\wedge\; H_1 \;\wedge\; H_2 \;\wedge\; \cdots \;\wedge\; H_n
\]
\vspace{0.1cm}
$H_i$: hipóteses auxiliares (instrumentos, condições, teorias subsidiárias\ldots)
\pause
\vspace{0.4cm}
\begin{alertblock}{Quando algo dá errado}
\[
\neg(T \wedge H_1 \wedge \cdots \wedge H_n) \equiv \neg T \vee \neg H_1 \vee \cdots \vee \neg H_n
\]
Sabemos que \impt{alguma} é falsa --- mas \impt{não qual}.
\end{alertblock}
\pause
\vspace{0.3cm}
\ntc{Resultado:} a falsificação \impt{nunca é conclusiva}.
\end{frame}

\begin{frame}{Duhem--Quine na prática: Urano vs.\ Mercúrio}
\ntc{Contexto:} No séc.~XIX, duas anomalias desafiaram a mecânica newtoniana. \\
As órbitas de \impt{Urano} e \impt{Mercúrio} não batiam com as predições.
\vspace{0.3cm}
\pause

Em cada caso, o teste envolvia:
\[
\text{Newton} \;\wedge\; \text{hipóteses auxiliares sobre o sistema solar}
\]
\pause
\vspace{0.2cm}
\ntc{Pergunta:} onde está o erro --- na teoria ou nas hipóteses auxiliares?
\pause
\vspace{0.3cm}
\begin{alertblock}{}
A tese de Duhem--Quine diz que a lógica \impt{não decide}. \\
A comunidade científica \impt{escolhe} onde localizar o erro.
\end{alertblock}
\end{frame}

\begin{frame}{Urano vs.\ Mercúrio --- duas respostas opostas}
\begin{columns}[t]
\column{0.48\textwidth}
\begin{block}{Urano (1846)}
Irregularidade na órbita.\\[0.15cm]
\ntc{Escolha:} revisar hipótese auxiliar \\ (``não há planetas além de Urano'').\\[0.15cm]
\impt{Netuno é postulado e descoberto.}\\[0.15cm]
Newton \emph{sobrevive}.
\end{block}

\column{0.48\textwidth}
\pause
\begin{block}{Mercúrio (1859--1915)}
Precessão anômala do periélio.\\[0.15cm]
\ntc{Mesma estratégia tentada:} postula-se \emph{Vulcano}, planeta entre Mercúrio e o Sol.\\[0.15cm]
\impt{Vulcano nunca é encontrado.}\\[0.15cm]
Relatividade Geral \emph{substitui} Newton.
\end{block}
\end{columns}
\pause
\vspace{0.4cm}
\begin{alertblock}{A moral}
Mesma teoria, mesma estratégia, \impt{desfechos opostos}. \\
Não há regra lógica dizendo quando persistir e quando abandonar.
\end{alertblock}
\end{frame}

\begin{frame}{3. Complexidade das situações reais}
Em experimentos reais, a rede de hipóteses auxiliares é \impt{vasta}:
\vspace{0.3cm}
\pause
\begin{itemize}
  \item Teorias dos instrumentos
  \item Correções estatísticas
  \item Suposições sobre condições de fundo
  \item Pressupostos matemáticos
\end{itemize}
\pause
\vspace{0.5cm}
\begin{alertblock}{}
Decidir o que foi ``falsificado'' exige \impt{julgamento científico}. \\
Não existe procedimento algorítmico.
\end{alertblock}
\end{frame}

\begin{frame}{4. Inadequação histórica --- Copérnico}
\begin{block}{A revolução copernicana}
Ao nascer, o heliocentrismo era:
\begin{itemize}
  \item \impt{Menos preciso} que Ptolomeu (aperfeiçoado)
  \item \impt{Contestado empiricamente} (sem paralaxe estelar)
  \item \impt{Incompatível com a física} aristotélica aceita (A Terra é feita de elemento terra cuja natureza é estar em repouso. Movimento exige causa atuante. Pedras caem em direção ao centro do universo.)
\end{itemize}
\end{block}
\pause
\vspace{0.5cm}
\ntc{Um falsificacionista rigoroso teria rejeitado Copérnico.}
\pause
\vspace{0.3cm}
\begin{alertblock}{}
A comunidade o manteve e desenvolveu --- \impt{e acertou}.
\end{alertblock}
\end{frame}


\begin{frame}{Balanço final}
O falsificacionismo é valioso como:
\vspace{0.2cm}
\begin{itemize}
  \item Antídoto contra dogmatismo
  \item Crítica ao indutivismo ingênuo
\end{itemize}
\pause
\vspace{0.4cm}
Mas falha em três frentes:
\vspace{0.2cm}
\pause
\begin{itemize}
  \item \impt{Logicamente} --- Duhem--Quine impede falsificação estrita
  \pause
  \item \impt{Epistemologicamente} --- enunciados observacionais são falíveis
  \pause
  \item \impt{Historicamente} --- não descreve como a ciência de fato progride
\end{itemize}
\pause
\vspace{0.4cm}
\ntc{Próximos passos:} Lakatos (programas de pesquisa) e Kuhn (paradigmas).
\end{frame}

% =====================================================
\section{Síntese}
% =====================================================

\begin{frame}{Síntese da aula}
\begin{block}{Falsificacionismo ingênuo}
Falseabilidade como demarcação; teorias ousadas.
\end{block}
\pause
\vspace{0.2cm}
\begin{block}{Falsificacionismo sofisticado}
Sequências de teorias; novas predições; rejeição de modificações \emph{ad hoc}.
\end{block}
\pause
\vspace{0.2cm}
\begin{alertblock}{Limitações}
Observação carregada de teoria; Duhem--Quine; inadequação histórica.
\end{alertblock}
\pause
\vspace{0.4cm}
\ntc{Lição geral:} critérios lógicos, por si sós, não explicam a ciência. \\
A ciência é atividade \impt{humana, histórica e social}.
\end{frame}

\begin{frame}[allowframebreaks]{Referências}
\small
\begin{thebibliography}{9}
\bibitem[Chalmers(1993)]{chalmers1993}
Chalmers, A.\ F. (1993).
\newblock \emph{O que é ciência afinal?}.
\newblock São Paulo: Brasiliense.

\bibitem[Popper(1972)]{popper1972}
Popper, K.\ R. (1972).
\newblock \emph{A lógica da pesquisa científica}.
\newblock São Paulo: Cultrix.

\bibitem[Lakatos(1978)]{lakatos1978}
Lakatos, I. (1978).
\newblock \emph{The methodology of scientific research programmes}.
\newblock Cambridge: Cambridge University Press.
\end{thebibliography}
\end{frame}

\end{document}
